\chapter{Khung Mẫu Truyện Từ Đời Sống Học Đường Và Học Thêm}
\markboth{Khung Mẫu Truyện Từ Đời Sống Học Đường Và Học Thêm}{Khung Mẫu Truyện Từ Đời Sống Học Đường Và Học Thêm}

\begin{center}
	\textit{\color{bookprimary}``Không phải ai cũng bí ý. Nhiều người chỉ thiếu một cái khung đủ gọn để viết ra điều mình đã thấy.''}

	\textit{\color{bookprimary}``Not everyone lacks ideas. Many people simply lack a frame simple enough to write what they have already seen.''}
\end{center}

\ngancach

\begin{goigiaovien}
Phần này không kể thêm truyện hoàn chỉnh.
\textit{(This section does not tell more complete stories.)}

Nó đưa ra các khung mẫu rất ngắn để giáo viên, phụ huynh, học sinh hoặc người biên soạn có thể tự biến thành truyện song ngữ theo đúng tinh thần của cả cuốn sách.
\textit{(It provides short templates so teachers, parents, students, or editors can turn them into bilingual stories in the spirit of the whole book.)}
\end{goigiaovien}

\latcat{Khung 1. Đứa Trẻ Vừa Ăn Vừa Xem}{Mở cảnh bằng một bữa cơm có đủ người nhưng ai cũng cầm máy. Chọn một chi tiết rất thật như tiếng muỗng chạm bát, tiếng video lọt ra, hay ánh mắt người lớn vừa ăn vừa cúi. Sau đó đặt một câu hỏi ngắn nhưng đau: trong bữa này, có ai thật sự ở cùng ai không?}{Open with a family meal where everyone is physically present but each person holds a screen. Choose one very real detail such as the sound of utensils, a leaking video sound, or adults eating while looking down. Then end with a painful short question: in this meal, is anyone truly with anyone?}

\latcat{Khung 2. Đi Học Thêm Trong Cơn Buồn Ngủ}{Dùng hình ảnh một học sinh ngồi sau xe lúc tối, mắt díp lại sau buổi học thêm. Không cần lên giọng nhiều. Chỉ cần tả đúng cơ thể mệt, cặp sách nặng, và lời động viên quen thuộc của người lớn. Cú chốt nên xoáy vào việc mệt mỏi được khen như thể đó là phẩm chất.}{Use the image of a student on the back of a motorbike late at night, eyes closing after tutoring. Do not over-preach. Simply describe the exhausted body, the heavy schoolbag, and the familiar encouragement from adults. The closing punch should expose how exhaustion is praised as if it were virtue.}

\latcat{Khung 3. Trong Lớp Nhưng Tâm Trí Ngoài Màn Hình}{Mở bằng một tình huống rất nhỏ: giáo viên đang hỏi, học sinh ngẩng lên kịp nhưng đầu không có gì. Từ đó kể ra nhịp phân tán vô hình do điện thoại để lại ngay cả khi điện thoại không còn trên tay. Kết câu bằng nhận định rằng cái máy có thể đã bỏ xuống, nhưng sự đứt mạch vẫn còn nằm trong đầu.}{Begin with a small moment: the teacher asks a question, the student looks up in time but has nothing in mind. From there describe the invisible fragmentation left by phones even when they are no longer in the hand. End with the point that the device may be set down, but the broken mental rhythm still remains inside.}

\latcat{Khung 4. Bảng Điểm Đẹp, Khuôn Mặt Xấu}{Chọn một học sinh có kết quả ổn nhưng sắc mặt ngày càng nặng. Đặt cạnh nhau hai hình ảnh đối lập: lời khen thành tích và biểu hiện kiệt sức. Truyện sẽ mạnh nếu tránh kết luận ồn ào mà chỉ cho người đọc tự thấy sự lệch giữa cái được khoe và cái đang hỏng.}{Choose a student with decent scores but an increasingly heavy face. Place two opposite images side by side: praise for achievement and visible exhaustion. The story becomes stronger if it avoids loud conclusions and simply lets the reader see the mismatch between what is displayed and what is decaying.}

\latcat{Khung 5. Một Nhà Bốn Người, Bốn Màn Hình}{Dùng bối cảnh quen thuộc trong phòng khách. Không ai cãi nhau, không ai đập phá, không ai hỗn. Chính cái yên quá mức ấy mới đáng sợ. Kết luận nên nhấn rằng nhiều gia đình trông hòa thuận vì mỗi người đã rút vào một hang phát sáng riêng.}{Use a familiar living-room setting. No one argues, no one breaks anything, no one behaves badly. That very excessive quiet is what makes it disturbing. The ending should highlight that many families look harmonious only because each person has retreated into a private glowing cave.}

\latcat{Khung 6. Thêm Một Lớp Để Chữa Một Nỗi Lo}{Nhân vật trung tâm có thể là phụ huynh chứ không phải học sinh. Cứ mỗi lần lo, người lớn lại thêm lớp, thêm thầy, thêm bài. Truyện nên bóc ra rằng đôi khi thứ được chữa không phải lỗ hổng học tập của đứa trẻ mà là cơn bất an của người lớn.}{The central figure can be a parent rather than a student. Each time anxiety rises, the adult adds another class, teacher, or worksheet. The story should expose that what is being treated is not always the child's learning gap but the adult's anxiety.}

\latcat{Khung 7. Điện Thoại Làm Mất Khoảng Trống}{Chọn một khoảnh khắc ngắn như đứng chờ, ngồi trên xe, hay nghỉ giữa hai tiết. Trước kia đó có thể là lúc nghĩ ngợi hoặc nhìn đời một chút. Nay nó bị nuốt trọn bởi động tác rút máy rất nhanh. Cú chốt nên nói rằng khi con người mất sạch khoảng trống, họ cũng mất luôn nơi để ý nghĩ lớn lên.}{Pick a small waiting moment such as standing in line, riding as a passenger, or resting between periods. It once could have been a time to reflect or simply observe life. Now it is swallowed by the quick gesture of pulling out the phone. The closing line should say that when people lose all empty space, they also lose the place where thoughts grow.}

\latcat{Khung 8. Học Mẹo Nhiều, Nghĩ Ít Dần}{Lấy bối cảnh một lớp thêm nơi giáo viên truyền công thức giải nhanh. Đừng phủ nhận giá trị của mẹo hoàn toàn. Hãy chỉ ra rằng mẹo chỉ đẹp khi đi sau hiểu biết. Khi mẹo thế chỗ tư duy, học sinh sẽ nghiện đi đường tắt và sợ những câu hỏi cần hiểu thật.}{Set the scene in a tutoring class where the teacher passes along fast-solving tricks. Do not deny the value of shortcuts completely. Show that shortcuts are only beautiful when they come after understanding. When they replace thinking, students become addicted to shortcuts and afraid of questions requiring real understanding.}

\latcat{Khung 9. Một Cuộc Gọi Hữu Ích, Cả Ngày Bị Giật}{Đây là mẫu truyện công bằng. Mở đầu bằng việc điện thoại thật sự giúp một việc quan trọng. Sau đó xoay sang phần còn lại của ngày bị bẻ nhỏ bởi hàng chục thông báo vô nghĩa. Khung này hợp để viết kiểu không phiến diện: công nhận lợi ích thật nhưng vẫn chỉ ra giá phải trả.}{This is a balanced story template. Open with the phone genuinely helping in an important way. Then turn to the rest of the day being chopped apart by dozens of meaningless notifications. This frame works well when you want fairness: acknowledge a real benefit while still exposing the cost.}

\latcat{Khung 10. Người Lớn Mắng Điều Mình Đang Làm}{Chọn một cảnh cha mẹ hoặc giáo viên nhắc trẻ bỏ điện thoại, nhưng ngay sau đó chính người lớn lại cúi vào màn hình. Không cần chửi gắt. Chỉ cần đặt hai hành động cạnh nhau là đủ làm lộ ra sự bất nhất. Câu cuối nên nặng ở chỗ: trẻ con học bằng mắt rất nhanh.}{Choose a scene where a parent or teacher tells a child to put the phone away, only to look down at their own screen immediately afterward. No harsh scolding is needed. Simply placing the two actions side by side exposes the inconsistency. The last line should hit hard: children learn very quickly through their eyes.}

\latcat{Khung 11. Được Khen Chăm Chỉ Nhưng Không Còn Vui}{Mẫu này đi theo chiều tâm lý. Nhân vật làm rất nhiều việc đúng nhưng ngày càng không thấy vui, không còn hào hứng, và bắt đầu ghét những thứ từng thích. Đừng kể dài. Chỉ cần vài dấu hiệu đủ thật để người đọc nhận ra một kiểu kiệt sức âm thầm.}{This template follows the psychological route. The character does many things correctly but feels less joy over time, loses enthusiasm, and starts disliking things they once loved. Do not make it long. A few real signs are enough to reveal a quiet form of burnout.}

\latcat{Khung 12. Cắt Bớt Một Thứ, Sống Lại Một Chút}{Mẫu cuối không chỉ phê phán. Hãy kể một thay đổi nhỏ nhưng thật: bớt một lớp thêm, bớt một giờ màn hình, ngủ đủ hơn, ăn tử tế hơn, hay ngồi chữa lỗi sai đến nơi đến chốn. Câu kết nên nhẹ hơn các mẫu khác, như một nhắc nhở rằng con người không nhất thiết phải thắng bằng cách làm mình nát ra.}{The final template should not only criticize. Tell one small but real change: one fewer tutoring class, one less hour of screen time, better sleep, proper meals, or patiently correcting mistakes. The ending should be gentler than the others, reminding readers that human beings do not have to destroy themselves in order to improve.}

\begin{tuhocnhanh}
1. Trong 12 khung trên, khung nào giống đời sống quanh em nhất? \textit{(Which of the 12 frames most resembles the life around you?)}

2. Nếu phải tự viết một truyện ngắn 1 trang, em sẽ chọn chi tiết thật nào để mở đầu? \textit{(If you had to write a one-page story, which real-life detail would you choose to open with?)}

3. Muốn truyện bớt sáo, em cần bỏ bớt kiểu kết luận nào? \textit{(If you want the story to avoid cliché, what kind of conclusion should you remove?)}
\end{tuhocnhanh}

\begin{tongketchuong}
Một cuốn sách kiểu này mạnh không chỉ vì luận điểm nặng.
\textit{(A book like this is strong not only because its arguments are heavy.)}

Nó mạnh vì nhặt đúng những cảnh đời nhỏ mà ai cũng thấy nhưng ít ai chịu gọi tên cho ra nhẽ.
\textit{(It is strong because it picks the right small scenes everyone has seen but few are willing to name clearly.)}
\end{tongketchuong}

\ngancach